\SourceLeaf{378}%
\needspace{8\baselineskip}
\begin{center}[2.]\end{center}

Die allgemeine Auflösung der Aufgabe, eine krumme Fläche in den kleinsten Theilen ähnlich auf eine Ebene zu entwerfen, scheint folgende zu sein:

Es sei $V=0$ die Gleichung der krummen Fläche und
\[
dV=p\,dx+q\,dy+r\,dz.
\]
Man integrire die Gleichung
\[
q\,dx-p\,dy+\BareRoot{-1}\,.\,\BareRoot{(pp+qq+rr)}\,.\,dz=0,
\]
indem man sie mit $V=0$ verbindet. Das Integral sei
\[
W=\mathrm{Const.}
\]
Eben so sei das Integral von
\[
q\,dx-p\,dy-\BareRoot{-1}\,.\,\BareRoot{(pp+qq+rr)}\,.\,dz=0:
\]
\[
W'=\mathrm{Const.}
\]
Alsdann ist die vollständige Auflösung:
\[
\begin{aligned}
\xi+\BareRoot{-1}\,.\,\eta&=f(W)\\
\xi-\BareRoot{-1}\,.\,\eta&=f'(W),
\end{aligned}
\]
wo $f,f'$ willkürliche Functionen bedeuten; oder vielmehr die Integration von
\[
q\,dx-p\,dy+\BareRoot{-1}\,.\,\BareRoot{(pp+qq+rr)}\,.\,dz=0
\]
gebe
\[
T+U\BareRoot{-1}=\mathrm{Const.},
\]
sodann ist
\[
\xi\pm\BareRoot{-1}\,.\,\eta=\mathrm{funct.}\,(T\pm U\BareRoot{-1}).
\]

Es sei das Integral von
\[
dx^2+dy^2+dz^2=0
\]
mit $V=0$ verbunden:
\[
T+U\BareRoot{-1}=\mathrm{Const.},
\]
so findet das obige Resultat statt.
